%%%%%%%%%% source scan idx 94 | folio 87 %%%%%%%%%%

\medskip
1. \underline{集合层}

消灭圈层将定义为第~3 节对任意集合层定义的左正合函子 $\Psi$ 的导出函子的值.

\medskip
1.1. \underline{Galois 层.}

1.1.1. 设 $Y$ 为域 $k$ 上的概形, 其可分闭包为 $\bar{k}$, 置 $\bar{Y}=Y\otimes_k\bar{k}$. 设 $G$ 为预有限群, 并设
\begin{equation*}
u:G\longrightarrow\operatorname{Gal}(\bar{k}/k)
\end{equation*}
为从 $G$ 到 $\bar{k}$ 在 $k$ 上的 Galois 群的连续同态. Galois 群 $\operatorname{Gal}(\bar{k}/k)$ 通过结构转移左作用于 $\bar{Y}$. 群 $G$ 通过 $u$ 作用于 $\bar{Y}$.

设 $\mathcal F$ 为 $\bar Y$（的平展景）上的集合层. $G$ 在 $\mathcal F$ 上的一个与 $G$ 在 $\bar Y$ 上的作用\underline{相容}的作用, 是 $G$ 通过对 $(\bar Y,\mathcal F)$ 的自同构所给出的作用, 并且在 $\bar Y$ 上诱导给定的作用. 换言之, 它是一个满足 $\sigma(gh)=\sigma(g)\,\sigma(h)$ 的同构系统
\begin{equation*}
\sigma(g):u(g)_*\mathcal F\longrightarrow\mathcal F
\end{equation*}
若 $G$ 作用于 $\mathcal F$, 则对于 $Y$ 上的任意平展 $U$, 令 $\bar U=U\otimes_k\bar k$ 为 $\bar Y$ 上的平展概形, 群 $G$ 作用于 $\mathcal F(\bar U)$.

\underline{定义 1.1.2.} \underline{如果对 $Y$ 上的每个拟紧平展 $U$, $G$ 在离散集合 $\mathcal F(\bar U)$ 上的作用都连续, 则称 $G$ 连续作用于 $\mathcal F$.}

只考虑仿射 $U$ 也得到相同的定义.

若 $\mathcal F$ 是 $Y$ 上的集合层, 其在 $\bar Y$ 上的逆像为 $\bar{\mathcal F}$, 则 $\operatorname{Gal}(\bar k/k)$ 通过结构转移作用于 $\bar{\mathcal F}$, 且与其在 $\bar Y$ 上的作用相容.

\underline{回顾 1.1.3.}
\begin{enumerate}[label=(\roman*),leftmargin=2.7em]
\item \underline{$\operatorname{Gal}(\bar k/k)$ 在 $\bar{\mathcal F}$ 上的作用连续.}

%%%%%%%%%% source scan idx 95 | folio 88 %%%%%%%%%%

\item \underline{函子}
\begin{equation*}
\mathcal F\longmapsto\bar{\mathcal F}
\end{equation*}
\underline{其中 $\bar{\mathcal F}$ 带有 $\operatorname{Gal}(\bar k/k)$ 的作用, 给出从 $Y$ 上集合层所成范畴到 $\bar Y$ 上带有 $\operatorname{Gal}(\bar k/k)$ 连续作用且与 $\operatorname{Gal}(\bar k/k)$ 在 $\bar Y$ 上的作用相容的集合层所成范畴的等价.}
\end{enumerate}

将 $\bar k$ 写为有限 Galois 扩张 $k_i$ 的正向极限. 设
\begin{equation*}
\varphi_i:Y_i=Y\otimes_k k_i\longrightarrow Y,
\qquad
\bar\varphi:\bar Y\longrightarrow Y
\end{equation*}
为投影. 对 $Y$ 上的仿射平展 $U$, 置 $U_i=U\otimes_k k_i$. 则
\begin{equation*}
\bar{\mathcal F}(\bar U)
=\varinjlim_i(\varphi_i^*\mathcal F)(U_i).
\end{equation*}
在 $(\varphi_i^*\mathcal F)(U_i)$ 上, 群 $\operatorname{Gal}(\bar k/k)$ 通过 $\operatorname{Gal}(k_i/k)$ 作用, 这证明了 (i).

(ii) 中函子的右伴随为函子
\begin{equation*}
\mathcal G\longmapsto
(\bar\varphi_*\mathcal G)^{\operatorname{Gal}(\bar k/k)}.
\end{equation*}
对 $Y$ 上的每个 $\mathcal F$, 有
\begin{equation*}
\bar\varphi_*\bar\varphi^*\mathcal F
=\varinjlim_i\varphi_{i*}\varphi_i^*\mathcal F.
\end{equation*}
因此, 为证明
\begin{equation*}
\mathcal F\xrightarrow{\sim}
(\bar\varphi_*\bar\varphi^*\mathcal F)^{
  \operatorname{Gal}(\bar k/k)},
\end{equation*}
只须对每个 $i$ 证明
\begin{equation*}
\mathcal F\xrightarrow{\sim}
(\varphi_{i*}\varphi_i^*\mathcal F)^{
  \operatorname{Gal}(k_i/k)}.
\end{equation*}
甚至只须证明这个箭头在从 $k$ 到 $k_i$ 的标量扩张后成为同构; 此时 $\operatorname{Spec}(k_i)$ 成为若干个 $\operatorname{Spec}(k)$ 副本的和, $\operatorname{Gal}(k_i/k)$ 在这些副本上单纯传递地置换, 于是断言显然成立.

因此 (ii) 中的正合函子是全忠实的. 最后, 只须验证箭头
\begin{equation*}
\bar\varphi^*(\bar\varphi_*\mathcal G)^{
  \operatorname{Gal}(\bar k/k)}\longrightarrow\mathcal G
\end{equation*}
为满态射. 设 $\bar\varphi_i:\bar Y\longrightarrow Y_i$. 由假设可知
\begin{equation*}
\varinjlim_i
\bar\varphi_i^*(\bar\varphi_{i*}\mathcal G)^{
  \operatorname{Gal}(\bar k/k_i)}
\longrightarrow\mathcal G
\end{equation*}
%%%%%%%%%% source scan idx 96 | folio 89 %%%%%%%%%%

为满态射.

我们有
\begin{equation*}
\bar\varphi^*(\bar\varphi_*\mathcal G)^{
  \operatorname{Gal}(\bar k/k)}
=
\bar\varphi_i^*\!\left(
  \varphi_i^*\varphi_{i*}
  \left((\bar\varphi_{i*}\mathcal G)^{
    \operatorname{Gal}(\bar k/k_i)}\right)^{
    \operatorname{Gal}(k_i/k)}
\right),
\end{equation*}
再用与上述相同的 Galois 下降论证, 注意到对 $Y_i$ 上带有 $\operatorname{Gal}(k_i/k)$ 作用且与其在 $Y_i$ 上的作用相容的 $\mathcal G'$, 有
\begin{equation*}
\varphi_i^*(\varphi_{i*}\mathcal G')^{
  \operatorname{Gal}(k_i/k)}
\xrightarrow{\sim}\mathcal G'.
\end{equation*}
即得结论.
